\section{Conclusão e Trabalhos Futuros}

Este trabalho apresentou um estudo empírico e comparativo de seis algoritmos de aprendizado de máquina aplicados ao \textit{Adult Income Dataset}, com o objetivo de prever a faixa de renda de indivíduos com base em características sociodemográficas. A adoção de uma metodologia robusta, fundamentada na separação estrita de dados para prevenir o vazamento de informações e na validação cruzada em dez partições (\textit{10-fold cross-validation}), garantiu que as estimativas de desempenho refletissem a real capacidade de generalização dos modelos.

Os resultados evidenciaram que a \textbf{Árvore de Decisão} consistiu na arquitetura mais adequada para este cenário específico, alcançando o melhor equilíbrio preditivo com um F1-Score de 0.6603. O sucesso deste estimador ressalta a importância de utilizar algoritmos capazes de realizar partições espaciais não lineares e de lidar nativamente com a heterogeneidade das variáveis categóricas e numéricas presentes em dados censitários. Em contrapartida, o fraco desempenho do Naive Bayes Gaussiano ilustrou claramente os riscos práticos de se aplicar modelos probabilísticos estritos quando as premissas matemáticas subjacentes, como a independência condicional e a normalidade das distribuições contínuas, são fortemente violadas pelo conjunto de dados.

Ademais, observou-se que modelos sensíveis à métrica de distância e à otimização por gradiente (KNN, Regressão Logística, Linear SVM e Rede Neural MLP) apresentaram desempenhos altamente competitivos e aglomerados. Esse fenômeno comprova a eficácia da etapa de pré-processamento elaborada no \textit{pipeline}, especialmente a padronização \textit{Z-score}, a qual permitiu a convergência estável e eficiente destas abordagens no espaço multidimensional de características.

Como propostas para trabalhos futuros, visando expandir a fronteira de resultados alcançada por esta pesquisa, sugere-se a investigação de métodos baseados em comitês (\textit{Ensembles}), tais como \textit{Random Forest} ou \textit{Gradient Boosting}. Tais arquiteturas tendem a mitigar a variância e o risco de sobreajuste (\textit{overfitting}) de árvores individuais, elevando a acurácia global. Recomenda-se também a exploração de técnicas de amostragem sintética ao nível dos dados, como o algoritmo SMOTE, para um tratamento mais incisivo do desbalanceamento de classes, bem como a adoção da otimização Bayesiana para uma varredura mais refinada e computacionalmente eficiente dos hiperparâmetros.
\clearpage